% ==============================================================================
% 03_ids_ips.tex — UE SIS589
% ==============================================================================
\chapter{Systèmes de Détection et de Prévention d'Intrusion}
\label{chap:ids}

\epigraph{\itshape « Un IDS sans analyste est une alarme sans gardien. »}
         {--- Proverbe SOC}

\begin{boxobjectifs}
\begin{enumerate}
  \item Distinguer IDS, IPS, NIDS, HIDS, EDR, XDR et NDR selon leur positionnement
        et leur mode de détection.
  \item Décrire les deux paradigmes de détection~: signatures vs anomalies.
  \item Lire et interpréter les alertes Snort/Suricata~; écrire une règle simple.
  \item Comprendre le fonctionnement d'un HIDS basé sur FIM et détection
        comportementale (Wazuh).
  \item Positionner les capteurs IDS dans une architecture réseau segmentée.
  \item Identifier les limites et techniques d'évasion des IDS classiques.
\end{enumerate}
\end{boxobjectifs}

% ==============================================================================
\section{Taxonomie des systèmes de détection}
\label{sec:ids-taxo}
% ==============================================================================

\subsection{Positionnement dans la défense en profondeur}

Les systèmes de détection constituent la deuxième ligne de défense~: là
où les contrôles préventifs (firewall, antivirus, hardening) cherchent à
\textbf{empêcher}, les IDS/IPS cherchent à \textbf{détecter} ce qui a
franchi ou contourné ces barrières.

\begin{table}[H]
\centering\small
\caption{Taxonomie des systèmes de détection et de prévention d'intrusion}
\label{tab:ids-taxo}
\rowcolors{2}{TableauPair}{TableauImpair}
\begin{tabular}{p{1.8cm}p{5cm}p{3.5cm}p{3.5cm}}
\toprule
\tableauHeader{Sigle} & \tableauHeader{Définition} &
\tableauHeader{Positionnement} & \tableauHeader{Action} \\
\midrule
\textbf{NIDS} & \textit{Network IDS} --- analyse le trafic réseau &
En coupure (inline) ou en écoute (\textit{span/tap}) & Alerte uniquement \\
\textbf{NIPS} & \textit{Network IPS} --- variante préventive du NIDS &
En coupure obligatoire & Alerte + blocage actif \\
\textbf{HIDS} & \textit{Host-based IDS} --- agent sur le système &
Sur chaque hôte surveillé & Alerte (FIM, logs, comportement) \\
\textbf{HIPS} & \textit{Host-based IPS} --- variante active &
Sur chaque hôte & Alerte + isolation processus \\
\textbf{EDR}  & \textit{Endpoint Detection \& Response} &
Agent hôte + console cloud & Détection + réponse automatisée \\
\textbf{NDR}  & \textit{Network Detection \& Response} &
Analyse flux réseau (NetFlow, PCAP) & Détection comportementale réseau \\
\textbf{XDR}  & \textit{Extended Detection \& Response} &
Corrélation HIDS + NIDS + SIEM + CTI & Réponse unifiée multi-couche \\
\bottomrule
\end{tabular}
\end{table}

\subsection{Mode passif vs mode inline}

\begin{description}
  \item[\termecle{Mode passif (tap/span)}] Le capteur reçoit une copie du
        trafic via un port miroir (SPAN) du switch ou un TAP optique.
        \textbf{Zéro impact sur la latence réseau.} Peut seulement alerter,
        pas bloquer. Mode recommandé pour un premier déploiement.
  \item[\termecle{Mode inline (IPS)}] Le capteur est inséré physiquement
        dans le flux réseau~: tout paquet le traverse.
        \textbf{Peut bloquer en temps réel.} Introduit une latence et un
        point de défaillance unique~: une panne de l'IPS coupe le réseau.
        Nécessite une configuration de \textit{bypass} matériel.
\end{description}

% ==============================================================================
\section{Paradigmes de détection}
\label{sec:ids-paradigmes}
% ==============================================================================

\subsection{Détection par signatures}

\begin{boxdef}
\textbf{Détection par signatures~:} Comparaison du trafic ou du comportement
observé avec une base de règles décrivant des patterns d'attaques connus.
Analogie~: l'antivirus à base de hash. Efficace pour les menaces connues~;
aveugle aux menaces inconnues (\textit{zero-day}) et aux variantes légèrement
modifiées.
\end{boxdef}

Avantages~: faible taux de faux positifs, déterminisme, explicabilité.
Inconvénients~: inefficace contre les zero-days, nécessite une mise à jour
permanente des signatures.

\subsection{Détection par anomalies}

\begin{boxdef}
\textbf{Détection par anomalies (comportementale)~:} Établissement d'une
\textit{baseline} du comportement normal (volume de trafic, ports utilisés,
horaires de connexion, patterns applicatifs), puis détection de toute déviation
statistiquement significative. Peut détecter des attaques inconnues~; génère
davantage de faux positifs pendant la phase d'apprentissage.
\end{boxdef}

Avantages~: capable de détecter des menaces inconnues et des attaques lentes
(\textit{slow and low}).
Inconvénients~: faux positifs élevés, opacité des modèles ML, phase
d'apprentissage longue (2--4~semaines typiquement).

\begin{boiteRemarque}{Complémentarité signatures et anomalies}
Les SIEM et EDR modernes combinent les deux approches~: signatures pour la
détection rapide et déterministe des menaces connues~; anomalies (souvent
via ML/UEBA) pour la détection de comportements déviants de comptes
compromis ou d'insiders malveillants. Cette complémentarité réduit les angles
morts de chaque approche individuelle.
\end{boiteRemarque}

% ==============================================================================
\section{Snort / Suricata~: NIDS open-source}
\label{sec:snort-suricata}
% ==============================================================================

\subsection{Présentation}

\textbf{Snort} (Cisco, 1998) et \textbf{Suricata} (OISF, 2010) sont les deux
NIDS/NIPS open-source de référence. Suricata est multi-thread et supporte
nativement le protocole EVE-JSON pour l'intégration SIEM. Les deux utilisent
une syntaxe de règles compatible et partagent les mêmes jeux de règles
communautaires (Emerging Threats, Snort Community Rules).

\subsection{Anatomie d'une règle Suricata}

\begin{lstlisting}[style=StyleCode,
  caption={Anatomie d'une règle Suricata},
  label={lst:suricata-rule}]
# Syntaxe :
# action proto src_ip src_port direction dst_ip dst_port (options)

alert tcp any any -> $HOME_NET 22 (
  msg:"SCAN SSH Brute Force Attempt";
  flow:to_server;
  threshold:type threshold, track by_src, count 5, seconds 60;
  classtype:attempted-admin;
  sid:2001219;
  rev:12;
  metadata:affected_product SSH, attack_target Server,
            deployment Perimeter, signature_severity Major,
            mitre_attack T1110.001;
)
\end{lstlisting}

\begin{table}[H]
\centering\small
\caption{Champs d'une règle Suricata}
\label{tab:suricata-champs}
\rowcolors{2}{TableauPair}{TableauImpair}
\begin{tabular}{p{2.5cm}p{4cm}p{7cm}}
\toprule
\tableauHeader{Champ} & \tableauHeader{Valeurs possibles} &
\tableauHeader{Description} \\
\midrule
\texttt{action}   & alert, pass, drop, reject &
Que faire quand la règle matche \\
\texttt{proto}    & tcp, udp, icmp, http, dns, tls &
Protocole réseau ou applicatif (Suricata supporte L7) \\
\texttt{src/dst}  & IP, CIDR, variables (\$HOME\_NET) &
Filtrage source et destination \\
\texttt{msg}      & Texte libre &
Description affichée dans l'alerte \\
\texttt{content}  & Chaîne ou hex &
Contenu recherché dans le payload \\
\texttt{sid}      & Entier unique &
Identifiant unique de la règle \\
\texttt{classtype}& attempted-admin, etc. &
Catégorie de l'alerte \\
\bottomrule
\end{tabular}
\end{table}

\begin{lstlisting}[style=StyleCode,
  caption={Règles Suricata utiles pour LiZbiZ-SIS},
  label={lst:suricata-rules-lizbiz}]
# ─── Détection scan de ports SYN ────────────────────────────────────────────
alert tcp any any -> $HOME_NET any (
  msg:"SCAN SYN Portscan Detected";
  flags:S,12; threshold:type both, track by_src, count 20, seconds 3;
  classtype:network-scan; sid:9000001; rev:1;)

# ─── SQLi dans requêtes HTTP (payload) ─────────────────────────────────────
alert http $EXTERNAL_NET any -> $HTTP_SERVERS $HTTP_PORTS (
  msg:"SQLI Union Select Detected";
  flow:established,to_server;
  http.uri; content:"UNION"; nocase; content:"SELECT"; nocase; distance:0;
  classtype:web-application-attack; sid:9000002; rev:1;)

# ─── Connexion vers nœud TOR (exit node connu) ──────────────────────────────
alert ip $HOME_NET any -> $TOR_EXIT_NODES any (
  msg:"POLICY Connexion vers noeud TOR sortant";
  classtype:policy-violation; sid:9000003; rev:1;)

# ─── Détection Cobalt Strike beacon (JA3 hash connu) ───────────────────────
alert tls any any -> any any (
  msg:"MALWARE Cobalt Strike Beacon TLS";
  ja3.hash; content:"a0e9f5d64349fb13191bc781f81f42e1";
  classtype:trojan-activity; sid:9000004; rev:1;)
\end{lstlisting}

\subsection{Format d'alerte EVE-JSON}

Suricata produit les alertes au format EVE-JSON, directement ingérable
par le SIEM~:

\begin{lstlisting}[style=StyleTerminal,
  caption={Alerte EVE-JSON Suricata},
  label={lst:eve-json}]
{
  "timestamp"     : "2026-04-12T03:14:22.000000+0000",
  "flow_id"       : 1234567890,
  "in_iface"      : "em0",
  "event_type"    : "alert",
  "src_ip"        : "203.0.113.42",
  "src_port"      : 54321,
  "dest_ip"       : "10.0.1.10",
  "dest_port"     : 22,
  "proto"         : "TCP",
  "alert"         : {
    "action"      : "allowed",
    "gid"         : 1,
    "signature_id": 9000001,
    "rev"         : 1,
    "signature"   : "SCAN SYN Portscan Detected",
    "category"    : "Network Scan",
    "severity"    : 2
  },
  "flow"          : { "pkts_toserver": 23, "pkts_toclient": 0 }
}
\end{lstlisting}

% ==============================================================================
\section{HIDS~: détection sur l'hôte}
\label{sec:hids}
% ==============================================================================

\subsection{Principes}

Un HIDS analyse des données locales à l'hôte~:
\begin{itemize}
  \item \textbf{FIM (\textit{File Integrity Monitoring})~:} calcul
        périodique de hash sur les fichiers système critiques~;
        alerte si une modification non autorisée est détectée.
  \item \textbf{Surveillance des journaux~:} analyse des logs locaux
        en temps réel (auth.log, Event Viewer).
  \item \textbf{Détection de rootkits~:} recherche de processus cachés,
        fichiers invisibles, hooks dans le noyau.
  \item \textbf{Conformité~:} vérification des configurations vs benchmarks
        CIS, PCI-DSS, ISO~27001.
\end{itemize}

\subsection{Wazuh~: HIDS open-source de référence}

\begin{boxdef}
\textbf{Wazuh~:} Plateforme open-source combinant HIDS (agent), SIEM léger,
FIM, détection de conformité et intégration avec MITRE ATT\&CK. Architecture
Manager/Agent~: l'agent collecte et transmet~; le manager corrèle et alerte.
Compatible Linux, Windows, macOS, Docker, Kubernetes.
\end{boxdef}

\begin{lstlisting}[style=StyleCode, language=YAML,
  caption={Configuration Wazuh --- FIM et détection rootkits},
  label={lst:wazuh-config}]
<!-- /var/ossec/etc/ossec.conf (extrait) -->
<ossec_config>

  <!-- FIM : surveiller les fichiers système critiques -->
  <syscheck>
    <frequency>3600</frequency>  <!-- verif toutes les heures -->
    <directories check_all="yes" report_changes="yes" realtime="yes">
      /etc,/usr/bin,/usr/sbin,/bin,/sbin
    </directories>
    <!-- Windows : fichiers systeme et registre -->
    <directories check_all="yes">%WINDIR%/System32</directories>
    <registry>HKEY_LOCAL_MACHINE\Software\Microsoft\Windows\CurrentVersion\Run</registry>
    <!-- Exclusions légitimes -->
    <ignore>/etc/mtab</ignore>
    <ignore>/etc/hosts.deny</ignore>
  </syscheck>

  <!-- Détection de rootkits -->
  <rootcheck>
    <frequency>43200</frequency>  <!-- 2 fois / jour -->
    <check_files>yes</check_files>
    <check_trojans>yes</check_trojans>
    <check_dev>yes</check_dev>
    <check_sys>yes</check_sys>
    <check_pids>yes</check_pids>
    <check_ports>yes</check_ports>
    <check_if>yes</check_if>
  </rootcheck>

  <!-- Commandes personnalisées (audit de comptes) -->
  <localfile>
    <log_format>command</log_format>
    <command>last -n 20</command>
    <frequency>360</frequency>
  </localfile>

</ossec_config>
\end{lstlisting}

% ==============================================================================
\section{EDR~: détection avancée sur le poste}
\label{sec:edr}
% ==============================================================================

\begin{boxdef}
\textbf{EDR (\textit{Endpoint Detection and Response})~:} Agent installé sur
chaque poste/serveur qui capture en continu la télémétrie comportementale
(processus, connexions réseau, accès fichiers, modifications de registre,
appels système), la transmet à une console centralisée pour corrélation,
et offre des capacités de réponse distante~: isolation réseau, kill de
processus, collecte de fichiers suspects.
\end{boxdef}

\begin{table}[H]
\centering\small
\caption{Comparaison antivirus classique vs EDR}
\label{tab:av-vs-edr}
\rowcolors{2}{TableauPair}{TableauImpair}
\begin{tabular}{p{4cm}p{5cm}p{5cm}}
\toprule
\tableauHeader{Critère} & \tableauHeader{Antivirus classique} &
\tableauHeader{EDR} \\
\midrule
Mode de détection & Signatures (hashes, patterns) &
Signatures + comportement + ML \\
Télémétrie & Événements fichiers uniquement &
Processus, réseau, mémoire, registre \\
Malwares \textit{fileless} & Non détecté &
Détecté (injection mémoire, LOLBins) \\
Réponse & Quarantaine fichier & Isolation réseau, kill process, forensics \\
Threat hunting & Non applicable & Requêtes sur l'historique comportemental \\
Déploiement & Agent léger & Agent avec overhead variable (5--15\% CPU) \\
\bottomrule
\end{tabular}
\end{table}

Solutions de référence~: \textbf{CrowdStrike Falcon}, \textbf{SentinelOne},
\textbf{Microsoft Defender for Endpoint}, \textbf{Elastic EDR} (open-source
partiel), \textbf{Velociraptor} (open-source, investigation forensique).

% ==============================================================================
\section{Positionnement des capteurs dans l'architecture réseau}
\label{sec:ids-positionnement}
% ==============================================================================

Le positionnement des capteurs IDS/IPS détermine leur visibilité et leur
efficacité. Pour une architecture segmentée en VLAN (cas LiZbiZ-SIS)~:

\begin{figure}[H]
\centering
\begin{tikzpicture}[
  zone/.style={rectangle,draw=CouleurPrimaire,fill=CouleurDef,
               rounded corners,minimum width=5cm,minimum height=0.8cm,
               font=\small\bfseries},
  ids/.style={rectangle,draw=CouleurBordAlerte,fill=CouleurAlerte,
              rounded corners,minimum width=2.8cm,minimum height=0.7cm,
              font=\small\bfseries},
  arrow/.style={-Stealth,thick,color=CouleurSecondaire}]

\node[zone] (inet)  at (0, 5)   {Internet};
\node[ids]  (nids1) at (3.5, 5) {NIDS\#1\\(périmètre)};
\node[zone] (fw)    at (0, 3.5) {Firewall pfSense};
\node[zone] (dmz)   at (0, 2)   {VLAN~10~: DMZ\\(Apache, Nginx)};
\node[ids]  (nids2) at (3.5, 2) {NIDS\#2\\(DMZ $\leftrightarrow$ LAN)};
\node[zone] (app)   at (0, 0.5) {VLAN~20~: Applicatif\\(ERP, MySQL)};
\node[zone] (usr)   at (0,-1)   {VLAN~30~: Utilisateurs\\(AD, postes)};
\node[ids]  (hids)  at (3.5,-0.3) {HIDS (Wazuh)\\sur chaque hôte};

\draw[arrow] (inet) -- (fw);
\draw[arrow] (fw)   -- (dmz);
\draw[arrow] (dmz)  -- (app);
\draw[arrow] (app)  -- (usr);

\draw[dashed,color=CouleurAccent] (nids1) -- (0.7,5);
\draw[dashed,color=CouleurAccent] (nids2) -- (0.7,2);
\draw[dashed,color=CouleurBordForensic] (hids) -- (0.7,0.5);
\draw[dashed,color=CouleurBordForensic] (hids.south) -- (0.7,-1);

\node[font=\tiny\itshape,color=CouleurBordAlerte] at (5.5,4.6)
  {Trafic N/S entrant};
\node[font=\tiny\itshape,color=CouleurBordAlerte] at (5.5,1.6)
  {Trafic DMZ$\to$LAN};
\node[font=\tiny\itshape,color=CouleurBordForensic] at (5.5,-0.3)
  {Télémétrie hôte};
\end{tikzpicture}
\caption{Positionnement des capteurs IDS dans l'architecture LiZbiZ-SIS}
\label{fig:ids-positionnement}
\end{figure}

\begin{itemize}
  \item \textbf{NIDS\#1 (périmètre)~:} Capture le trafic entre Internet
        et le firewall. Visibilité sur les attaques entrantes brutes.
  \item \textbf{NIDS\#2 (inter-VLAN)~:} Capture les flux entre la DMZ
        et les VLANs internes. Détecte les pivots et mouvements latéraux.
  \item \textbf{HIDS (Wazuh)~:} Déployé sur chaque serveur. Détecte
        les attaques qui ont franchi le réseau (modifications locales,
        exécutions suspectes, FIM).
\end{itemize}

% ==============================================================================
\section{Limites et techniques d'évasion des IDS}
\label{sec:ids-evasion}
% ==============================================================================

La connaissance des limites des IDS est indispensable pour les configurer
correctement et interpréter leurs silences.

\begin{table}[H]
\centering\small
\caption{Techniques d'évasion des IDS et contre-mesures}
\label{tab:ids-evasion}
\rowcolors{2}{TableauPair}{TableauImpair}
\begin{tabular}{p{3.5cm}p{5cm}p{5cm}}
\toprule
\tableauHeader{Technique (MITRE)} & \tableauHeader{Description} &
\tableauHeader{Contre-mesure} \\
\midrule
Fragmentation IP (T1001.002) &
Découpage du payload en fragments pour éviter les patterns &
Réassemblage IP avant inspection (IPS inline) \\
Encodage (T1027) &
Base64, XOR, Unicode pour masquer le contenu &
Décodage applicatif L7 dans l'IDS \\
Chiffrement TLS (T1573) &
Payload chiffré, signatures inefficaces &
Inspection TLS (TLS MITM), analyse JA3/JA3S \\
Trafic lent (\textit{slow scan}) &
Une sonde par heure, en dessous des seuils &
Fenêtres temporelles longues, UEBA \\
Living off the Land (T1218) &
Outils légitimes (powershell, certutil) &
Surveillance comportementale, liste blanche \\
Obfuscation réseau &
Tunneling DNS, ICMP, HTTPS &
Inspection protocolaire, analyse entropie \\
\bottomrule
\end{tabular}
\end{table}

\begin{boiteRemarque}{TLS et inspection}
Plus de 90\% du trafic Internet est aujourd'hui chiffré (TLS~1.3).
Un NIDS classique ne voit que les métadonnées (IP, ports, durée, volume)
et l'empreinte TLS (JA3 hash), pas le contenu. L'inspection TLS profonde
nécessite un proxy d'interception (MITM) avec implications légales et
de confidentialité. L'analyse JA3 permet de fingerprinter le client
et de détecter certains outils offensifs connus.
\end{boiteRemarque}

% ==============================================================================
\section*{Résumé}
% ==============================================================================

\begin{boxresume}
\begin{itemize}
  \item NIDS (réseau), HIDS (hôte), EDR (endpoint avancé), NDR (réseau
        comportemental), XDR (corrélation unifiée).
  \item Deux paradigmes~: signatures (menaces connues, peu de FP) et
        anomalies (menaces inconnues, FP élevés pendant l'apprentissage).
  \item Snort/Suricata~: syntaxe action/proto/src/dst/(options). Format
        EVE-JSON pour l'intégration SIEM.
  \item HIDS Wazuh~: FIM, surveillance logs, rootkits, conformité CIS.
  \item EDR~: télémétrie comportementale complète, détection
        \textit{fileless}, réponse distante.
  \item Positionnement~: périmètre (N/S), inter-VLAN (E/O), agent hôte.
  \item Limites~: fragmentation, chiffrement TLS, trafic lent, LOLBins.
\end{itemize}
\end{boxresume}

\begin{boxexo}
\textbf{Ex.~3.1 — Écriture de règle Suricata.}
Rédigez une règle Suricata détectant une tentative de \textit{reverse shell}
Netcat (\texttt{nc -e /bin/bash}) vers un port non standard (hors 80/443/22)
depuis le VLAN~20. Spécifiez action, mapping MITRE et SID unique.

\medskip\textbf{Ex.~3.2 — Analyse d'alerte EVE-JSON.}
À partir de l'alerte ci-dessous (fournie en TP), identifiez~:
(a)~l'attaque probable~; (b)~le TTP MITRE~; (c)~si le trafic est bloqué
ou seulement alerté~; (d)~les 3~étapes de qualification à suivre en N1.

\medskip\textbf{Ex.~3.3 — Plan de déploiement IDS.}
Proposez un plan de déploiement IDS/HIDS complet pour LiZbiZ-SIS en
précisant~: localisation de chaque capteur, mode (tap ou inline), type
(NIDS, HIDS, EDR), règles prioritaires à activer et métriques de suivi.
\end{boxexo}
